\chapter{UI guide}
\label{ch:guida_ui_en}

Tour of the SvelteKit tabs: dashboard, anagrafica (teachers /
classes / students / classrooms / \emph{ore}), assignments, monitor,
optimize (with iterative-diversified and branch-and-price modes),
diagnostics, runs.

\section*{Working hours configuration (Tab Ore)}

The \texttt{/ore} tab lets the school administrator configure the
working week:
\begin{itemize}
  \item which days are active (default: Mon-Sat);
  \item how many timetable slots each day has and at what times
        (default: 6 slots, 08:00-14:00).
\end{itemize}

Slots are stored as 0-based indices internally
(\texttt{WorkingHourSlot.slot\_index}), with explicit
\texttt{start\_time}/\texttt{end\_time} (HH:MM strings) used only
for display in the weekly calendar view and reports. Days carry a
\texttt{position} that the engine consumes as \texttt{day\_idx}, plus
a \texttt{legacy\_day\_number} (1..7) that keeps backward
compatibility with the existing \texttt{teacher\_unavailability} /
\texttt{class\_unavailability} / \texttt{classroom\_unavailability}
rows keyed on the historical DAYS=[1..6] / HOURS=[8..13]
convention.

The engine reads this configuration through
\texttt{engine/working\_hours\_config.py}, which falls back to the
legacy hardcoded defaults whenever the \texttt{working\_days} table
is missing or empty so old DBs keep working unchanged.

The unavailability matrices on the Teachers / Classes / Classrooms
tabs were migrated to a new reusable
\texttt{<WeeklyCalendarView />} component that consumes the same
configuration: identical color palette, click-cycle and keyboard
shortcuts (\texttt{H}/\texttt{P}/\texttt{E}/\texttt{D}/\texttt{A}/
\texttt{N} + click for HARD/PREFERRED/ENFORCED/SOFT/ALLOWED/RESET)
as the previous \texttt{AvailabilityMatrix}, but with columns
driven by the active working days and slot times rendered from
the configuration.

\medskip
\noindent
\textit{This chapter is an English summary of the corresponding Italian chapter. The full text remains in the Italian edition (\texttt{docs/manual/chapters/guida\_ui.tex}). For the implementation references see the engine modules in \texttt{engine/} and the API documentation at \texttt{docs/api.md}.}
